\documentclass{article}
\usepackage{graphicx} % Required for inserting images
\usepackage{amsfonts}

\title{lec422}
\author{Dylan Tang}
\date{April 2026}

\begin{document}

\maketitle

\section{Supervised Learning}
Fundamentally, we aim to optimize weights that can give us predictions $\hat{y}$
\[
    \hat{y} = X\hat{w}
\]
where $w \in \mathbb{R}^{n \times 1}$.
\\
\\
Recall the OLS problem, which minimizes squared loss:

Then, OLS regression gives us:
\[
    \arg\min_w |y - Xw|_2^2 
\]

\[
    \hat{w} = (X^TX)^{-1}X^Ty
\]

Also recall the ridge regression problem, which minimizes squared loss and an L2 regularizer:
\[
    \arg\min_w |y - Xw|_2^2  + |\lambda|^2_2
\]

\[
    \hat{w} = (X^TX + \lambda I)^{-1}X^Ty
\]

\textbf{Gradient Descent}:
\begin{itemize}
    \item Taking the inverse to get the optimum is extremely computationally intensive, and not feasible for large $n$.
    \item Instead, we can use gradient descent:
    \[
        w^{(t+1)} = w^{(t)} - \eta\nabla(\sum_{i = 1}^n (y- x^Tw)^T(y-x^Tw) + \lambda^T\lambda)
    \]
\end{itemize}

\textbf{Inductive Model}: Given $X, y$ as inputs, outputs $f(\theta)$
\begin{itemize}
    \item Linear models are inductive models. Given $(x_i, y_i)$ pairs, produces $f(\theta) = \hat{y} = \hat{\beta}_1X + \hat{\beta}_0 = X\hat{\beta}$
    \item Then, given a new data point $x_{new}$, we can find $\hat{y}_{new}$
    \[
        \hat{y}_{new} = x_{new}^T\hat{\beta}
    \]
\end{itemize}

\textbf{Transductive Model}: Given $X, y$ as inputs, outputs a data structure that organizes the inputs into $(X,y)$ pairs
\begin{itemize}
    \item K-nearest-neighbors (KNN), K-means clustering are transductive models
\end{itemize}

\subsection{Supervised Learning using Scikit}

\textbf{Step 1: Train-test Split}:
\begin{itemize}
    \item Always do train-test split before supervised learning.
\end{itemize}
\begin{verbatim}
    from sklearn.model_selection import train_test_split

    X_train, X_test, y_train, y_test = train_test_split(
        X, y, test_size=0.2, random_state=42
    )
\end{verbatim}

\textbf{Step 2: Fit the model}
\begin{verbatim}
    from sklearn.model_selection import train_test_split
    from sklearn.linear_model import LinearRegression
    
    X_train, X_test, y_train, y_test = train_test_split(
        X, y, test_size=0.2, random_state=42
    )
    
    model = LinearRegression()
    model.fit(X_train, y_train)
\end{verbatim}

\textbf{Step 3: Find errors on test set}
\begin{verbatim}
    from sklearn.model_selection import train_test_split
    from sklearn.linear_model import LinearRegression
    import numpy as np
    
    X_train, X_test, y_train, y_test = train_test_split(
        X, y, test_size=0.2, random_state=42
    )
    
    model = LinearRegression()
    model.fit(X_train, y_train)
    y_pred = model.predict(X_test)
    err = np.linalg.norm(y_test - y_pred)
\end{verbatim}

\section{Smoothing}
\begin{itemize}
    \item Recall $R^2 = \frac{SSE}{SST}$.
        \begin{itemize}
            \item $SST = \sum_{i = 1}^n(y_i - \bar{y})^2$ can be thought of as the error of the smoothest possible predictor $\bar{y}$.
        \end{itemize}
    \item Imagine several decision boundaries classifying labels $\{1 \dots n\}$.
    \begin{itemize}
        \item If those decision boundaries are spread out far apart, then it is a smooth model - less data dependent.
        \item  If those decision boundaries are very close together, then it is not smooth - more data dependent
    \end{itemize}
  
\end{itemize}

\end{document}
